\section*{Chapter \ref{ch:m2:inflation-markets} --- Inflation Markets}

\begin{solution}{exo:m2:inflation-markets:1}
$292.296 + \tfrac{14}{31}\times(296.311 - 292.296) = 294.10922$: August has 31
days and the 15th is day 15, so fourteen thirty-firsts of the way.
\end{solution}

\begin{solution}{exo:m2:inflation-markets:2}
Exactly $1.0267/1.0014 - 1 = 2.5265\%$; the difference, 2.53\%, overstates it
by a fraction of a basis point at these levels.
\end{solution}

\begin{solution}{exo:m2:inflation-markets:3}
The TIPS repays 100, its \omterm{def:m2:inflation-markets:linker}{deflation floor}; the gilt repays 97, since its
principal follows the index down. Coupons paid along the way were on the
adjusted principal in both cases.
\end{solution}

\begin{solution}{exo:m2:inflation-markets:4}
$100 \times (345/300 - 1.025^5) = \text{USD}~1.86$ million to the fund. None at the
forward index, $300 \times 1.025^5 = 339.42$.
\end{solution}

\begin{solution}{exo:m2:inflation-markets:5}
$294.10922/277.20274 = 1.06099$. On 15 July the \omterm{def:m2:inflation-markets:linker}{reference index} interpolated
between the April and May indices, lower than May and June while prices were
rising: the ratio was 1.04814.
\end{solution}

\begin{solution}{exo:m2:inflation-markets:6}
$-0.144 - 0.359 - 0.349 = -0.85$ percentage points over the three months, about
3.4 points a year at an annual rate. A breakeven read from such a bond without
adjustment looks 3.4 points too low.
\end{solution}

\begin{solution}{exo:m2:inflation-markets:7}
On 2010--2019 March is still the strongest month (0.355) and November the
weakest ($-0.376$); no month moves by more than 0.11 percentage points from the
2010--2024 estimate, the largest shift being June's. The pattern is stable; its
size is estimated to about a tenth of a point.
\end{solution}

\begin{solution}{exo:m2:inflation-markets:8}
The 9\% is inflation over the past year; the breakeven is the price of the
next ten years' average inflation, with risk and liquidity premia in it. The
position earns realised inflation against the breakeven only over the months
it is held, and in the meantime it loses or gains as the breakeven moves; the
high prints already published are in the price, through the known accrual.
The trade is a view that the breakeven will rise, or that realised inflation
over the holding period will beat what was priced, not a collection.
\end{solution}

\begin{solution}{pb:m2:inflation-markets:1}
\textbf{1.} 292.296 (the May index) and 296.311 (June).
\textbf{2.} 1.05445, the ratio of 292.296 to 277.20274.
\textbf{3.} 1.374\%, the growth from 292.296 to 296.311.
\textbf{4.} On 13 July, when the June index was published: over August the
\omterm{def:m2:inflation-markets:linker}{reference index} moves from the May to the June index, both then known.
\textbf{5.} 294.10922.
\textbf{6.} A real price of 99.865 (dirty) per 100 of adjusted principal; market
value $100 \times 0.99865 \times 1.05445$, USD~105.30 million.
\textbf{7.} $105.30 \times 1.374\%$, USD~1\,446\,444.
\textbf{8.} USD~12\,513 (0.14\% a year for 31 days), plus a cross term of USD~172.
\textbf{9.} $105.30 \times 2.30\% \times 31/360$, USD~208\,558.
\textbf{10.} USD~1\,250\,572.
\textbf{11.} USD~27\,366: the nominal bond's yield barely covers repo.
\textbf{12.} 17.8\% a year, against a breakeven of 2.53\%: the carry reflects a
single month's print, not the market's expectation for ten years.
\textbf{13.} $296.276/296.311 - 1 = -0.012\%$; carry $-202\,160$ dollars, the
repo cost with nothing to offset it.
\textbf{14.} USD~98\,921 per basis point; 12.6 basis points of \omterm{def:m2:inflation-markets:breakeven}{real yield} would
wipe out the month's carry.
\textbf{15.} The accrual is the same per dollar of value at any maturity, while
the real-yield risk grows with duration: a two-year TIPS earns the same
August accrual with a fifth of the price risk.
\textbf{16.} Not by itself: the forward price for the end of August includes the
known accrual, so a buyer on 29 July pays for it. It is profit only against
a position that did not own it, or if the \omterm{def:m2:inflation-markets:breakeven}{real yield} falls.
\textbf{17.} \omterm{def:m2:inflation-markets:breakeven}{Real yields} rose by two and a half points, and a real duration above
nine years turned that into a price loss larger than the year's indexation.
\textbf{18.} Very little: the index ratio was 1.054, so the index would have to fall
about 5\% by 2032 before the floor paid anything.
\textbf{19.} Named result: \emph{the \omterm{def:m2:inflation-markets:linker}{linker}'s carry} of USD~100 million face of the
ten-year TIPS over August 2022 was about USD~1.25 million, known on 13 July.
\textbf{20.} A \omterm{def:m2:inflation-markets:linker}{linker} protects its holder against the price level, not against
the real interest rate.
\end{solution}

\begin{solution}{iq:m2:inflation-markets:1}
The constant inflation rate at which a nominal bond and an inflation-linked
bond of the same maturity return the same, roughly the nominal yield minus the
\omterm{def:m2:inflation-markets:breakeven}{real yield}. It is a price: it includes a premium for bearing inflation risk,
minus the premium the less liquid \omterm{def:m2:inflation-markets:linker}{linker} must pay, plus convexity effects, and it
moves with flights to liquidity, as in March 2020.
\iqlookfor{the Fisher relation, and the premia that separate it from a forecast.}
\end{solution}

\begin{solution}{iq:m2:inflation-markets:2}
The principal is multiplied by the index ratio, the \omterm{def:m2:inflation-markets:linker}{reference index} (the
non-seasonally-adjusted CPI-U with a three-month lag, interpolated daily) over
its value on the dated date; coupons are the fixed real rate on that principal.
At maturity the holder receives the larger of the adjusted and the original
principal: deflation over the whole life is floored, though coupons paid
along the way on a lower principal are not.
\iqlookfor{the lagged, interpolated index, and the floor on principal only.}
\end{solution}

\begin{solution}{iq:m2:inflation-markets:3}
Because the principal grows over a month by the index change of two and three
months earlier, so a large print becomes a large, known accrual a few weeks
later; financed at the \omterm{def:m2:repo-and-specials:rate}{repo rate}, the position carries the accrual minus
repo. The carry goes to whoever held the bond before it was priced in; a buyer
after the print pays for it in the forward price.
\iqlookfor{the lag mechanics and the fact that known carry is priced.}
\end{solution}

\begin{solution}{iq:m2:inflation-markets:4}
Different instruments and costs: the swap has no funding and little balance
sheet, the bond must be financed and uses balance sheet; their liquidity
differs; the bond's \omterm{def:m2:inflation-markets:linker}{deflation floor} has value the swap lacks; supply and demand
differ (pension-fund receiving in swaps, Treasury issuance in bonds); the
lag conventions and the exact dates of the indices may not match; and the
bond's breakeven is quoted on a yield, the swap on a zero-coupon rate.
\iqlookfor{funding and balance sheet, the floor, flows, and conventions.}
\end{solution}

\begin{solution}{iq:m2:inflation-markets:5}
Estimate monthly factors from the history of the unadjusted index (for example,
each month's change less its year's mean, averaged over complete years, or the
official seasonal factors), check their stability over sub-periods, normalise
them to sum to zero over the year, and apply them to the forward index curve
built from swaps, so that forward monthly changes carry the calendar pattern;
then price the short \omterm{def:m2:inflation-markets:linker}{linkers} on that curve.
\iqlookfor{estimation, stability, normalisation, and applying it to forwards.}
\end{solution}

\begin{solution}{iq:m2:inflation-markets:6}
A conventions table per bond (index, lag, interpolation, rounding, dated date,
floor, coupon dates), with eight-month-lag gilts as a separate rule; an index
store keyed by month, loaded from each statistics office on its release day,
with checks against the previous value, its release calendar, and versioning
for rebasings; a pure function from (bond, date, index store) to the ratio,
tested on the issuers' published examples and ratio tables; an output that is
reconciled each morning against the issuers' own published ratios, and alerts
when an index is late or a bond's rule is missing.
\iqlookfor{conventions as data, a versioned index store, and reconciliation with
the issuers' tables.}
\end{solution}
